\documentclass[12pt,a4paper]{report}
\usepackage{graphicx}
\usepackage{amsmath}
\usepackage{geometry}
\usepackage{booktabs}
\usepackage{float}
\usepackage{titlesec}  
\usepackage{tocloft}   
\usepackage{setspace}
\usepackage{url}

\geometry{margin=2.5cm}

\begin{document}
\onehalfspacing

\begin{titlepage}
\centering
\vspace*{2cm}
\Huge\textbf{In-body Antigen Detection via Antibodies on a Transistor}\\[1cm]
\large CMPE 492 Midterm Report\\[2cm]
\textbf{Student:} Duran Kaan Altın \\
\textbf{Advisors:} Tuna Tuğcu and Kutlu Ülgen\\
\textbf{Date:} 31 March 2026
\end{titlepage}

\tableofcontents
\newpage

\chapter{Introduction}

\section{Broad Impact}

Cancer develops when normal regulatory mechanisms for cell division fail, allowing malignant cells to invade other tissues. This unchecked growth allows malignant cells to spread to distant organs through blood or lymphatic circulation, disrupting normal physiological functions~\cite{janeway2001}. Among the numerous biomarkers employed in cancer diagnostics, one of the most significant for certain breast cancer subtypes is the \textbf{Human Epidermal Growth Factor Receptor 2 (HER2)}. This receptor, located on the cell membrane, normally regulates cell proliferation and repair. However, when overexpressed, HER2 drives abnormal cell signaling that results in excessive cell division and tumor expansion~\cite{yarden2001}.

In HER2-positive cases, excessive receptor expression accelerates cell proliferation and tumor growth. In normal cells, HER2 forms heterodimers with other members of the epidermal growth factor receptor family, such as HER3, to mediate controlled cell division. In contrast, mutations in the \textit{ERBB2} gene, which encodes HER2, lead to its amplification and constitutive activation. This genetic alteration converts \textit{ERBB2} into an oncogene, continuously signaling the cell to proliferate, thereby promoting rapid and invasive tumor growth~\cite{kabel2017}.

Clinically, HER2 expression is commonly determined through immunoassays that rely on specific antibody binding, including \textbf{enzyme linked immunosorbent assay (ELISA)} and \textbf{immunohistochemistry (IHC)}. Although these methods are sensitive and well validated, they require labeled reagents, multiple preparation steps, and sophisticated laboratory instrumentation~\cite{harpaz2017,janeway2001ch9}. As a result, they are unsuitable for continuous, in body, or real time monitoring applications.

This research aims to address those limitations by investigating a new biosensing approach based on \textbf{graphene field effect transistors (G-FETs)}. The concept is to detect HER2 molecules by observing the electrical response caused when antibody functionalized graphene surfaces bind with the target antigen. The binding event induces a change in surface charge, which alters the channel conductance of the transistor producing a measurable electrical signal~\cite{bergveld1972,pena2018}. This fusion of immunology and nanoelectronics enables label free detection and the potential for real time biosensing.

Through comprehensive multiphysics simulations, the study explores how various parameters such as charge carrier mobility, analyte concentration, surface potential, and channel geometry influence sensor performance. These simulations aim to determine the optimal conditions for achieving maximum sensitivity, selectivity, and detection limit. 

At its core, this work seeks to demonstrate how a nanoscale electronic device can convert invisible molecular recognition events into quantifiable electrical signals, paving the way for \textbf{fast, noninvasive, and continuous in-body cancer diagnostics}.


\begin{figure}[H]
    \centering
    \includegraphics[width=0.8\linewidth]{image.png}
    \caption{Healthy Breast Cells~\cite{biodigital}.}
    \label{fig:placeholder}
\end{figure}

HER2-positive breast cancer is an aggressive form of breast cancer characterized by high levels of a protein present in certain cells called human epidermal growth factor receptor, or HER2.
In healthy breast cells, the cell surface proteins HER3 and HER2 join together, or dimerize. The connection of these two proteins signals for the cell to duplicate itself normally~\cite{biodigital}.

\begin{figure}[H]
    \centering
    \includegraphics[width=0.8\linewidth]{image1.png}
    \caption{Malignant Cell~\cite{biodigital}.}
    \label{fig:placeholder}
\end{figure}

During HER2-positive breast cancer, the gene ERBB2 incorrectly instructs the cell to create increased levels of HER2 protein on the cell's surface. The abnormal gene is known as an oncogene because it causes tumor growth~\cite{biodigital}.

\begin{figure}[H]
    \centering
    \includegraphics[width=0.8\linewidth]{image2.png}
    \caption{Uncontrolled Reproduction of Cancerous Cells~\cite{biodigital}.}
    \label{fig:placeholder}
\end{figure}

The protein HER2 is involved in cell growth and differentiation. When overexpressed or amplified, HER2 causes uncontrolled proliferation of malignant (cancerous) cells. The cancer cells replicate rapidly to form a breast cancer tumor~\cite{biodigital}.


\section{Ethical Considerations}

Although this project is based entirely on simulation and does not involve in-vivo or clinical testing, the ethical implications of potential downstream deployment should be acknowledged. If the proposed technology progresses toward clinical or implantable use, it must satisfy established medical device safety and biocompatibility requirements, including the evaluation framework defined in \textbf{ISO 10993}~\cite{iso10993}.


In addition, methodological transparency is essential. All modeling assumptions, parameter choices, and numerical procedures should be explicitly documented and justified to support reproducibility and responsible translation to future clinical studies.





\chapter{Project Definition and Planning}

\section{Problem Definition}

\begin{figure}[H]
    \centering
    \includegraphics[width=0.8\linewidth]{image3.png}
    \caption{Breast Anatomy~\cite{biodigital}.}
    \label{fig:placeholder}
\end{figure}

The breasts are mounds of fat and glandular tissue in the front of the chest. Men and women have similar breast tissue during childhood, but sex-specific hormones initiate female breast development during puberty. During adulthood, female breasts have a well-developed glandular structure that is responsible for milk production.

Breast cancer is caused by the uncontrolled growth of abnormal cells in the breast. While it can originate from any cell type, the most common forms of breast cancer arise from the lining cells of the lactiferous ducts or lobules.

The progression of breast cancer is divided into stages based on the size of the tumor, the involvement of lymph nodes, and the metastasis to other parts of the body~\cite{biodigital}.

\begin{figure}[H]
    \centering
    \includegraphics[width=0.8\linewidth]{image9.png}
    \caption{Stage 0 Breast Cancer~\cite{biodigital}.}
    \label{fig:placeholder}
\end{figure}

Stage 0 breast tumors are in situ; they remain in the area in which they originated and show no signs of spreading to nearby tissues~\cite{biodigital}.

\begin{figure}[H]
    \centering
    \includegraphics[width=0.8\linewidth]{image4.png}
    \caption{Stage I Breast Cancer~\cite{biodigital}.}
    \label{fig:placeholder}
\end{figure}

Stage I breast tumors are invasive and have started to spread to the fatty breast tissue or to nearby lymph nodes. Tumor growth is small or nonexistent~\cite{biodigital}.

\begin{figure}[H]
    \centering
    \includegraphics[width=0.8\linewidth]{image5.png}
    \caption{Stage II Breast Cancer~\cite{biodigital}.}
    \label{fig:placeholder}
\end{figure}

Stage II breast tumors are invasive. The tumor may have grown in size and spread to additional lymph nodes~\cite{biodigital}.

\begin{figure}[H]
    \centering
    \includegraphics[width=0.8\linewidth]{image6.png}
    \caption{Stage III Breast Cancer~\cite{biodigital}.}
    \label{fig:placeholder}
\end{figure}

Stage III breast tumors are invasive and have spread to the lymph nodes around the collarbone or those deep within the breast. The tumor may continue to grow in size and may have spread to the chest wall or skin~\cite{biodigital}.

\begin{figure}[H]
    \centering
    \includegraphics[width=0.8\linewidth]{image7.png}
    \caption{Stage IV Breast Cancer~\cite{biodigital}.}
    \label{fig:placeholder}
\end{figure}

Stage IV breast tumors are invasive and have spread beyond the localized area to the lungs, bones, liver, brain, or to distant lymph nodes~\cite{biodigital}.


Breast cancer typically progresses through well-characterized stages, beginning with localized abnormal proliferation and potentially advancing to invasive and metastatic disease. Early on, abnormal cells begin multiplying where they shouldn’t, usually inside ducts or lobules. If left unchecked, they start invading surrounding tissue, spreading to lymph nodes, and eventually reaching distant organs like the lungs or brain~\cite{janeway2001,weissleder2004}.

This stepwise spread is classified into \textbf{stages}, which help doctors determine how advanced the disease is and how to treat it:

\begin{itemize}
    \item \textbf{Stage I:} The tumor is small and may just be entering nearby fatty tissue or lymph nodes.
    \item \textbf{Stage II:} The cancer is more established; the tumor is larger or has spread to additional lymph nodes.
    \item \textbf{Stage III:} The spread is now serious, possibly reaching the collarbone area or the chest wall.
    \item \textbf{Stage IV:} This is full-blown metastasis. Cancer has reached distant organs like the liver, lungs, or bones.
\end{itemize}

At this stage, the discussion shifts to the technical considerations that motivate the proposed sensing approach.

A key prognostic and predictive biomarker in breast cancer is \textbf{HER2}, a transmembrane receptor that regulates signaling pathways associated with cell proliferation and survival. It functions as a transmembrane receptor that acts as a regulatory control element for cell division. In some cancers, the HER2 gene becomes overactive, flooding the cell surface with HER2 proteins. This hijacks the cell’s growth signals, leading to runaway cell division and resistance to treatment~\cite{yarden2001}. 

Clinically, about 15-30 percent of breast cancer cases show \textbf{HER2 overexpression}, and that usually means a worse prognosis~\cite{kabel2017}. For this reason, monitoring HER2 expression is not only diagnostic but also critical for guiding treatment aggressiveness and evaluating therapeutic response.
 

HER2 expression is commonly assessed using \textbf{ELISA} and \textbf{IHC}. While these assays are clinically validated, they are not designed for continuous monitoring or in-body, real-time sensing due to sample preparation requirements and laboratory instrumentation dependencies~\cite{harpaz2017,janeway2001ch9}.


This limitation motivates the central problem addressed in this study.

A label-free electronic approach capable of transducing HER2 binding events into an electrical signal could enable faster and potentially continuous monitoring compared to conventional assays. One candidate platform for this purpose is the \textbf{graphene field-effect transistor (G-FET)}. Graphene exhibits high sensitivity to surface charge fluctuations, which makes it well-suited for conversion of molecular binding events into measurable electrical responses~\cite{bergveld1972,pena2018}.


When HER2 molecules bind to antibodies anchored on the graphene surface, they change the local surface charge. This causes a shift in the transistor’s current, a signal we can measure. The relationship looks something like this:

\begin{equation}
\Delta I_{DS} = \frac{W}{L} \, e \, \mu \, V_{DS} \, \Delta n \propto N
\end{equation}

Where:
\begin{itemize}
    \item $\Delta I_{DS}$ is the change in current,
    \item $N$ is the number of bound HER2 molecules,
    \item and the rest are physical constants or dimensions of the transistor.
\end{itemize}

However, accurately modeling this process presents significant challenges. Antigen transport by diffusion, electrostatic screening effects in ionic biological fluids, and charge transport within the graphene channel occur simultaneously and are strongly coupled. These phenomena are governed by nonlinear equations whose interactions must be considered to obtain physically realistic predictions.

Accordingly, the \textbf{goal of this study} is to develop a \textbf{multiphysics simulation framework} that integrates biochemical reaction kinetics, diffusion-driven transport, and electronic transduction within a unified computational model. This framework enables quantitative prediction of GFET biosensor response under varying physiological conditions and provides design insight for achieving reliable in-body operation.





\section{Solution Definition}
A COMSOL Multiphysics simulation model is suggested as a solution to examine HER2-antibody binding and how it affects the G-FET channel current. Transistor response under realistic flow and ionic conditions is predicted by combining biological kinetics with electronic transport. The process incorporates:
\begin{itemize}
  \item Reaction Engineering Module for antigen-antibody kinetics (0D)
  \item Transport of Diluted Species and Surface Reaction Modules for 3D diffusion and surface binding
  \item Semiconductor and AC/DC Modules for electrical current response (\( \Delta I_{ds} \))
\end{itemize}

In this term, the work is focused on "diffusion-only" transport in a lymph node inspired geometry, providing a transport limited baseline for later binding and GFET coupling. 

\section{Project Planning}
\begin{center}
\begin{tabular}{|l|c|p{7cm}|p{4cm}|}
\hline
\textbf{Work Package} & \textbf{Duration} & \textbf{Description} & \textbf{Deliverables} \\
\hline
WP1 & M1-M3& Literature review, CAD geometry, mesh setup & Fall Midterm Report\\
WP2 & M3-M4& Simulation of HER2 binding and flow transport & Fall Final Report\\
WP3 & M4-M7& Optimization of selectivity and sensitivity & Spring Midterm Report \\
WP4 & M7-M8& Reporting and visualization & Spring Final Report\\
\hline
\end{tabular}
\end{center}

\subsection{Time and Resource Estimation}
\begin{itemize}
  \item Software: COMSOL Multiphysics v5.4+
  \item Hardware: GPU-enabled workstation 
  \item Estimated simulation time: 8-12 hours per configuration
\end{itemize}

\subsection{Success Criteria}
\begin{itemize}
  \item Correlation between HER2 concentration and \( \Delta I_{ds} \)
  \item Achieve simulated limit of detection (LOD) $\leq 1$ pM
  \item Stable transistor response under both static and flow conditions
\end{itemize}

\subsection{Risk Analysis}
\begin{center}
\begin{tabular}{|p{4cm}|p{3cm}|p{6cm}|}
\hline
\textbf{Risk} & \textbf{Impact} & \textbf{Mitigation} \\
\hline
Incomplete HER2 data & Medium & Substitute with another cancer biomarker such as CA15-3 \\
COMSOL convergence issues & High & Mesh refinement and adaptive solver configuration \\
Hardware limitations & Medium & Utilize personal GPU server as backup\\
Unexpected nonlinear effects & Medium & Conduct parameter sweeps and recalibration \\
\hline
\end{tabular}
\end{center}

\chapter{Related Work}

The transduction of biological phenomena into electrical signals has been an active area of research for several decades. A foundational contribution was made in 1972 by \textbf{Bergveld}, who introduced the ion-sensitive field-effect transistor (ISFET). This work demonstrated that ionic variations in an aqueous environment could be measured using a solid-state electronic device, thereby establishing the conceptual basis for modern bioelectronic sensing technologies~\cite{bergveld1972}.

Subsequent research extended this principle by incorporating biological recognition elements into FET architectures. In a comprehensive review, \textbf{De Moraes and Kubota (2016)} systematically classified FET-based immunosensors according to their signal transduction mechanisms, including electrochemical, optical, and mechanical approaches~\cite{demoraes2016}. These developments highlighted the versatility of FET platforms for biosensing applications.

More recently, the emergence of graphene as a channel material has significantly advanced the field. Owing to its atomic-scale thickness and high charge carrier mobility, graphene exhibits exceptional sensitivity to surface charge perturbations. This property has been widely exploited in biosensor design. \textbf{Pe\~na-Bahamonde et al. (2018)} provided an extensive analysis of graphene-based biosensors and identified their high sensitivity and biocompatibility as key advantages for in-body diagnostic applications~\cite{pena2018}.

Several experimental demonstrations have validated the potential of graphene FET biosensors. \textbf{Seo et al. (2020)} reported a graphene-based FET capable of detecting SARS-CoV-2 antigens with sub-femtomolar sensitivity, illustrating both rapid response and high detection performance~\cite{seo2020}. In parallel, \textbf{Kim, Cho, and Yu (2018)} investigated the integration of carbon nanomaterials into flexible and wearable bioelectronic systems, ultimately expanding the applicability of biosensors beyond conventional laboratory environments~\cite{kim2018}.

In the context of cancer diagnostics, recent studies have demonstrated notable progress. \textbf{Xu et al. (2022)} developed an HER2 biosensor based on a functionalized graphene oxide FET, achieving detection limits below 10~pM with high selectivity, which represents a significant advancement for early-stage breast cancer detection~\cite{xu2022}. Similarly, \textbf{Zhang et al. (2021)} presented a multiplexed GFET platform capable of simultaneously detecting multiple cancer biomarkers through surface charge and electrostatic potential modulation~\cite{zhang2021}.

The present study builds directly upon this body of work while introducing a key methodological distinction. Rather than focusing on experimental fabrication or measurement, the emphasis is placed on computational modeling of HER2-sensitive G-FET behavior. By integrating experimentally informed parameters with reaction-diffusion transport and semiconductor physics models, this work establishes a predictive simulation framework that bridges theoretical analysis and practical biosensor design.


The ultimate objective is to quantitatively characterize how \textbf{antibody-antigen binding} modulates charge transport in the transistor channel, with sufficient predictive accuracy to inform the engineering design of future deployable, in-body cancer detection systems.



\chapter{Methodology}

\section{Understanding The Field Effect Transistor Operation}

Before introducing the simulation framework, it is necessary to clarify the physical quantity measured by a transistor-based biosensor. A \textbf{field-effect transistor (FET)} is a solid-state device in which the conductivity of a semiconductor channel is controlled by an electric field applied via a gate electrode. The device comprises three primary terminals: the \textit{source}, \textit{drain}, and \textit{gate}.

A gate-to-channel voltage modifies the electrostatic potential within the channel, resulting in a change in charge carrier density. This carrier modulation directly alters the drain-source current, allowing variations in surface potential or external electric fields to be converted into an electrical output signal.


In a \textbf{graphene field-effect transistor (G-FET)}, the channel is a single layer of graphene placed on an insulating layer (usually SiO$_2$) with an underlying doped silicon substrate acting as a back gate. Graphene’s high carrier mobility allows even small surface charge changes to cause measurable shifts in current. This property makes G-FETs particularly suitable for biochemical sensing applications.


\begin{figure}[H]
    \centering
    \includegraphics[width=0.8\linewidth]{image8.png}
    \caption{Schematic representation of GFET~\cite{seo2020}.}
    \label{fig:gfet_schematic}
\end{figure}

As shown in Figure~\ref{fig:gfet_schematic}, antibodies are immobilized on the graphene surface via linker molecules. When HER2 antigens bind to these antibodies, they alter the local surface potential. This charge perturbation modifies the carrier concentration in the graphene channel, leading to a measurable change in the drain current ($\Delta I_{DS}$). The current concentration relationship typically follows a saturation-type curve, where the response increases rapidly with concentration before reaching equilibrium.

The measured parameter in this biosensor is therefore the \textbf{drain–source current modulation}, expressed as:

\begin{equation}
\Delta I_{DS} = \frac{W}{L}\, e\, \mu\, V_{DS}\, \Delta n \propto N
\end{equation}

where $W$ and $L$ are the channel width and length, $e$ is the elementary charge, $\mu$ is the carrier mobility, $V_{DS}$ is the drain–source voltage, $\Delta n$ is the carrier density change, and $N$ represents the number of bound HER2 molecules.

This electrical transduction mechanism bridges biochemistry and electronics: molecular binding events occurring at the graphene surface are directly reflected as electrical signals in the transistor’s output. The upcoming sections model these phenomena step by step, starting from molecular kinetics and diffusion to the resulting electrical current response.

\section{Stage 1: 0D Kinetic Model}
Using COMSOL's Reaction Engineering interface, HER2-antibody binding is represented as:
\[
  r = k_f[Ag][Ab] - k_r[AgAb]
\]
where $k_f$ and $k_r$ are forward and reverse rate constants. The equilibrium constant $K_d$ is determined for integration into higher-dimensional simulations.

\section{Stage 2: 3D Transport and Surface Binding}
Transport of Diluted Species (TDS) and Surface Reaction interfaces simulate diffusion, convection, and surface binding:
\[
  \frac{\partial c_i}{\partial t} + \nabla \cdot (-D_i \nabla c_i) + u \cdot \nabla c_i = 0
\]
This setup allows investigation of HER2 concentration gradients, flow effects, and surface reaction kinetics.

\section{Stage 3: Electrical Response of GFET}
The transistor current modulation is calculated using:
\[
  \Delta I_{ds} = \frac{w}{l} e \mu V_{ds} \Delta n \propto N
\]
where $N$ is the number of bound HER2 molecules. Extracted metrics include transconductance ($g_m$) and charge-neutrality voltage ($V_{Dirac}$) shifts.

\section{Validation}
Simulation results will compare flow vs. static environments, varying HER2 concentrations (0.5-1000 pM), and ionic strengths to analyze Debye screening effects.

\section{Sensitivity and Detection Limit Estimation}

To evaluate the detectability of individual binding events in a compact sensing node,
the graphene FET (GFET) signal was modeled under realistic node-module constraints.
Starting from the relation

\begin{equation}
\Delta I_{ds} = 
\left( \frac{w}{l} \right) e \mu V_{ds} \Delta n , \qquad
\Delta n = \frac{\alpha N}{A_{\text{eff}}},
\end{equation}

the minimum detectable bound number is

\begin{equation}
N_{\min} =
\frac{\Delta I_{ds,\min} A_{\text{eff}}}
     {\alpha \left(\tfrac{w}{l}\right) e \mu V_{ds}} .
\end{equation}

\noindent
The parameters were set to match the constraints of an embedded node module:

\begin{itemize}
    \item Channel width \(w = 20~\mu\text{m}\), length \(l = 2~\mu\text{m}\) (\(w/l = 10\))
    \item Carrier mobility \(\mu = 0.1~\text{m}^2/\text{V·s}\) (\(\approx 1000~\text{cm}^2/\text{V·s}\))
    \item Drain bias \(V_{ds} = 50~\text{mV}\)
    \item Effective area \(A_{\text{eff}} \approx 4\times10^{-11}~\text{m}^2\)
    \item Coupling efficiency \(\alpha = 0.01\) (typical) – \(0.03\) (optimized)
    \item Readout noise floor \(\Delta I_{ds,\min} = 10{-}50~\text{pA}_{\text{RMS}}\)
\end{itemize}

\noindent
Substituting these values yields:

\begin{center}
\begin{tabular}{lcccc}
\toprule
$\Delta I_{ds,\min}$ (pA) & $\alpha$ & $N_{\min}$ & Interpretation \\ \midrule
50 & 0.01 & $\approx 25$ & typical node resolution \\
10 & 0.01 & $\approx 5$ & optimized node electronics \\
50 & 0.03 & $\approx 8$ & improved surface coupling \\
\bottomrule
\end{tabular}
\end{center}

Each bound HER2 molecule contributes approximately
\(\Delta I_{1} \approx 2~\text{pA}\times\alpha\),
so 5–25 molecules are needed to exceed realistic noise levels.
This corresponds to the previously predicted LOD range of \(0.5{-}1.0~\text{pM}\),
confirming that the simulated device performance is physically attainable in a
miniaturized biomedical node. Single-molecule readability is theoretically possible only
under ultra-low-noise or dry-state conditions.






\chapter{Requirements Specification}
\begin{center}
\begin{tabular}{|c|c|c|}
\hline
\textbf{Symbol} & \textbf{Meaning} & \textbf{Units} \\
\hline
C & HER2 concentration & pM \\
$\theta$ & Site occupancy & -- \\
$K_d$ & Dissociation constant & pM \\
$g_m$ & Transconductance & $\mu$A/V \\
$V_{Dirac}$ & Charge-neutrality voltage & V \\
$\Delta I_{ds}$ & Change in drain current & $\mu$A \\
LOD & Limit of detection & pM \\
\hline
\end{tabular}
\end{center}

\section*{Functional Requirements}
\begin{enumerate}
  \item Model HER2-antibody binding in both static and flow environments.
  \item Compute $\Delta I_{ds}$ as a function of HER2 concentration.
  \item Generate calibration plots and estimate LOD.
\end{enumerate}

\section*{Non-functional Requirements}
\begin{itemize}
  \item Simulation reproducibility within 5\% error.
  \item GPU computation time $<$ 12 hours per configuration.
  \item Export results in CSV and PNG formats for analysis.
\end{itemize}

\chapter{Design}

This chapter describes the conceptual design of the proposed sensing system and the modeling architecture used to support it. The focus is on how biological, physical, and electronic components are structured and how information flows between them, rather than on implementation-level details.

\section{Information Structure}

The proposed system is designed as a hierarchical sensing pipeline that transforms biochemical phenomena into electrical signals suitable for electronic readout. Unlike conventional software systems, the information structure in this project is governed by physical processes operating at multiple spatial scales.

At the highest level, the system consists of the following information layers:

\begin{itemize}
    \item \textbf{Transport layer:} Spatial and temporal distribution of HER2 antigens within tissue-scale volumes.
    \item \textbf{Surface interaction layer:} Antigen-antibody binding events at the functionalized graphene surface.
    \item \textbf{Transduction layer:} Conversion of surface charge perturbations into carrier density modulation in the GFET channel.
    \item \textbf{Electrical output layer:} Drain-source current changes measured by external readout circuitry.
\end{itemize}

Each layer produces well-defined physical outputs that serve as inputs to the next stage. The work completed in this term focuses on the transport layer, which defines the maximum achievable sensor input under realistic in-body conditions.

\section{Information Flow}

Information flow in the proposed system is dictated by physical transport mechanisms rather than digital communication. HER2 concentration acts as the primary information carrier and propagates through tissue via diffusion.

The diffusion-only simulations demonstrate that information flow is strongly influenced by spatial heterogeneity. Antigen exposure at the sensing region depends not only on inlet concentration but also on diffusivity contrasts between tissue compartments. As shown in the results, concentration information reaches the cortex earlier and with higher amplitude than the medulla, introducing a measurable delay.

This behavior establishes a fundamental system-level constraint: regardless of sensor sensitivity, detection cannot occur before sufficient antigen transport to the sensing interface. Consequently, transport dynamics play a critical role in determining sensor response time and must be considered as part of the overall system design.

\section{System Design}

The proposed system is designed as a modular biosensing architecture that integrates biological transport modeling with nanoelectronic signal transduction. The system can be decomposed into four interacting subsystems:

\begin{enumerate}
    \item \textbf{Biological transport subsystem:} Models antigen delivery through tissue using diffusion-dominated transport.
    \item \textbf{Biochemical interaction subsystem:} Describes antigen-antibody binding at the graphene surface using reaction kinetics.
    \item \textbf{Electronic transduction subsystem:} Represents the GFET device that converts surface charge variations into current modulation.
    \item \textbf{Signal acquisition subsystem:} Encompasses the electronic circuitry responsible for current measurement and noise suppression.
\end{enumerate}

In the scope of this term, the biological transport subsystem is fully implemented and validated through multiphysics simulation. The remaining subsystems are defined analytically and structurally, forming the basis for subsequent implementation in later project phases.

This staged design approach reduces system complexity while ensuring that each subsystem is physically and computationally justified before full integration.

\section{User Interface Design}

A user-facing interface is not applicable for the current stage of the project. The system is intended as an embedded biomedical sensing platform rather than an interactive software application. Visualization and analysis are performed using scientific post-processing tools to support engineering evaluation rather than end-user interaction.

\chapter{Implementation and Testing}

This chapter documents the implementation details and validation methodology for the diffusion-only transport model developed during this term.

\section{Implementation}

The transport model was implemented in COMSOL Multiphysics using the \textit{Transport of Diluted Species} interface. A three-dimensional lymph-node-inspired geometry was constructed and partitioned into cortex and medulla subdomains to enable heterogeneous material assignment.

Spatially varying diffusion coefficients were applied to each subdomain, and transient simulations were performed under fixed inlet concentration conditions. Boundary conditions were selected to reflect transport-limited in-body exposure scenarios while maintaining numerical stability over long simulation times (up to 6000~s).

Post-processing workflows were implemented to extract surface concentration distributions, internal concentration slices, diffusive flux streamlines, and volume-averaged concentration metrics for quantitative comparison between compartments.

\section{Testing}

Testing in this project focuses on numerical correctness and physical consistency rather than functional software testing. Validation steps included:

\begin{itemize}
    \item Verification of mass conservation under diffusion-only conditions.
    \item Sensitivity analysis with respect to diffusivity ratios between cortex and medulla.
    \item Mesh and time-step stability checks to ensure convergence of transient solutions.
\end{itemize}

Parametric sweeps over the diffusivity ratio
\begin{equation}
r = \frac{D_{\text{medulla}}}{D_{\text{cortex}}}
\end{equation}
were used to evaluate model robustness and to confirm that observed concentration delays are a direct consequence of transport heterogeneity rather than numerical artifacts.

\section{Deployment}

The current implementation is not deployed as an end-user system. Instead, it serves as a validated simulation backbone for future project phases. Planned deployment includes coupling the transport model with GFET electrostatic simulations and integrating the full sensing chain into a predictive framework for in-body biosensor design.



\chapter{Results and Discussion}

\section{Diffusion-Only Transport in a Lymph Node-Inspired Geometry}

This term, the multiphysics pipeline was intentionally narrowed to a diffusion-only transport model in order to establish a physically consistent and computationally workable baseline for \textit{in-body} antigen delivery. The key question addressed here is not yet the transistor electrostatics, but the more fundamental prerequisite: \textbf{Can HER2 molecules reach the sensing region fast enough and in sufficient quantity to be detectable?}

The lymph-node-inspired geometry was constructed as a simplified 3D organ-scale volume with (i) an outer capsule, (ii) an internal subdomain representing the medulla, and (iii) inlet/outlet openings representing lymphatic connections. Figure~\ref{fig:ln_geometry} shows the final geometry and domain partition used for the diffusion study.

\begin{figure}[H]
    \centering
    \includegraphics[width=0.85\linewidth]{model.png}
    \caption{Lymph-node-inspired 3D geometry used in the diffusion-only model. The outer domain represents the cortex/capsule region and the inner domain represents the medulla. Openings are included to support later flow and boundary-driven studies; in this term they are used as concentration boundaries.}
    \label{fig:ln_geometry}
\end{figure}



\subsection{Governing equation and boundary conditions}

Transport was modeled using Fickian diffusion with heterogeneous diffusivity:
\begin{equation}
\frac{\partial c}{\partial t} = \nabla \cdot \left(D(\mathbf{x}) \nabla c \right),
\end{equation}
where $c(\mathbf{x},t)$ is the antigen concentration and $D(\mathbf{x})$ is spatially piecewise constant:
\begin{equation}
D(\mathbf{x})=
\begin{cases}
D_{\text{cortex}}, & \mathbf{x}\in \Omega_{\text{cortex}}\\
D_{\text{medulla}}, & \mathbf{x}\in \Omega_{\text{medulla}}.
\end{cases}
\end{equation}

An inlet opening was assigned a fixed concentration $c=c_0$ (Dirichlet), representing an antigen-carrying lymphatic inflow. Remaining external boundaries were treated as no-flux:
\begin{equation}
\mathbf{n}\cdot\left(-D\nabla c\right)=0,
\end{equation}
and the initial condition was $c(\mathbf{x},0)=0$.

For consistency with biochemical units, the inlet concentration was defined as $c_0 = 10~\text{pM}$ and converted to SI units in COMSOL:
\begin{equation}
c_{0,\text{SI}} = c_0 \cdot 10^{-12}\, \frac{\text{mol}}{\text{L}} \cdot 1000\,\frac{\text{L}}{\text{m}^3}
= 10^{-8}\, \frac{\text{mol}}{\text{m}^3}.
\end{equation}

\section{Spatiotemporal Concentration Profiles}

Figure~\ref{fig:conc_surface} shows the concentration distribution on the lymph node surface at $t=200~\text{s}$, demonstrating that antigen exposure is highly localized around the inlet boundary at early times. A mid-plane slice visualization is given in Figure~\ref{fig:conc_slice}, which highlights delayed penetration into the interior subdomain (medulla) compared to the cortex region.

\begin{figure}[H]
    \centering
    \includegraphics[width=0.85\linewidth]{concsurface.png}
    \caption{Surface concentration distribution at $t=200~\text{s}$ for the diffusion-only model. Antigen concentration remains highest near the inlet boundary, indicating transport-limited exposure at early times.}
    \label{fig:conc_surface}
\end{figure}

\begin{figure}[H]
    \centering
    \includegraphics[width=0.85\linewidth]{slice.png}
    \caption{Slice visualization of concentration in the lymph-node-inspired geometry at $t=200~\text{s}$. The interior subdomain (medulla) exhibits delayed concentration build-up compared to the outer cortex region.}
    \label{fig:conc_slice}
\end{figure}

\section{Diffusive Flux Field Evolution}

Although the model does not include convection in this term, the diffusive flux field,
\begin{equation}
\mathbf{J} = -D\nabla c,
\end{equation}
provides an intuitive picture of transport pathways from the inlet region into the organ volume. Figure~\ref{fig:flux_streamlines} shows streamlines of total diffusive flux colored by concentration. The result reinforces that the dominant transport direction is driven by steep local gradients near the inlet and then spreads into the cortex before substantially reaching the medulla.

\begin{figure}[H]
    \centering
    \includegraphics[width=0.95\linewidth]{flux.png}
    \caption{Streamlines of total diffusive flux $\mathbf{J}=-D\nabla c$ colored by concentration at $t=200~\text{s}$ (example shown for $r=D_{\text{medulla}}/D_{\text{cortex}}=0.1$). The flux is strongest near the inlet boundary, where gradients are steepest.}
    \label{fig:flux_streamlines}
\end{figure}

To distinguish transient transport effects from intrinsic compartmental
limitations, the diffusive flux field was further examined at
$t = 6000~\text{s}$ for the same heterogeneous case
$r = D_{\text{medulla}}/D_{\text{cortex}} = 0.1$.
At this time scale, the system approaches a quasi-steady state, and
overall concentration gradients are significantly reduced compared to
early times.

Figure~\ref{fig:flux_streamlines_6000s} shows that although the magnitude
of the diffusive flux decreases, the dominant transport pathways remain
topologically unchanged.
Flux streamlines continue to preferentially traverse the cortex region,
while penetration into the medulla remains limited.
This indicates that reduced medullary exposure is not merely a transient
delay but a persistent consequence of heterogeneous diffusivity.

From a transport perspective, this result demonstrates that increasing
observation time alone does not eliminate compartmental shielding when
$D_{\text{medulla}} \ll D_{\text{cortex}}$.
Instead, the tissue-scale diffusivity contrast imposes a fundamental
constraint on antigen delivery to internal regions.

\begin{figure}[H]
    \centering
    \includegraphics[width=0.95\linewidth]{flux_6000s.png}
    \caption{Streamlines of total diffusive flux
    $\mathbf{J} = -D\nabla c$ colored by concentration at
    $t = 6000~\text{s}$ for $r = D_{\text{medulla}}/D_{\text{cortex}} = 0.1$.
    Although flux magnitudes are reduced relative to early times, transport
    pathways remain dominated by the cortex, indicating persistent
    transport limitation of the medulla at steady state.}
    \label{fig:flux_streamlines_6000s}
\end{figure}



\section{Cortex vs. Medulla Uptake Delay}

To quantify heterogeneity-induced delay, volume-average concentrations were computed separately in the cortex and medulla domains:
\begin{equation}
\langle c(t)\rangle_{\Omega} = \frac{1}{|\Omega|}\int_{\Omega} c(\mathbf{x},t)\, dV,
\end{equation}
for $\Omega=\Omega_{\text{cortex}}$ and $\Omega=\Omega_{\text{medulla}}$.

Figure~\ref{fig:cortex_medulla_curve} shows the resulting time-dependent uptake curves. The cortex rises rapidly due to proximity to the inlet boundary and larger effective diffusivity, while the medulla response is substantially delayed. This separation is not a numerical artifact; it is the expected outcome of heterogeneous transport and internal compartmentalization.

\begin{figure}[H]
    \centering
    \includegraphics[width=0.95\linewidth]{plot1.png}
    \caption{Volume-average concentration vs. time in cortex and medulla domains for the diffusion-only model. The medulla exhibits delayed uptake relative to the cortex, demonstrating transport-limited exposure in internal regions.}
    \label{fig:cortex_medulla_curve}
\end{figure}


A short parametric sweep was performed on the diffusivity ratio
\begin{equation}
r = \frac{D_{\text{medulla}}}{D_{\text{cortex}}},
\end{equation}
to capture how strongly internal transport is penalized when the medulla is more restrictive. Figure~\ref{fig:diffusivity_sweep} illustrates the sweep configuration and indicates that decreasing $r$ increases the cortex-medulla disparity in concentration buildup over the fixed observation window.

\begin{figure}[H]
    \centering
    \includegraphics[width=0.85\linewidth]{parametric_sweep.png}
    \caption{Parametric sweep configuration over the diffusivity ratio $r=D_{\text{medulla}}/D_{\text{cortex}}$. Lower $r$ corresponds to stronger transport limitation in the medulla and therefore larger uptake delay.}
    \label{fig:diffusivity_sweep}
\end{figure}

\section{Implications for GFET-Based HER2 Detection}

The diffusion-only results directly inform the feasibility of in-body GFET detection by clarifying a hard constraint: \textbf{the transistor cannot detect molecules that have not yet arrived at the sensing surface.} In practice, a GFET measures an electrical signal that depends on the local near-surface concentration $c_s(t)$, not the bulk-average concentration in the full organ.

In the intended biosensor chain:
\begin{enumerate}
    \item \textbf{Transport sets the sensor input:} diffusion delivers antigen to the sensing region with a time-dependent exposure $c_s(t)$.
    \item \textbf{Binding converts concentration to occupancy:} assuming Langmuir-type binding at the functionalized graphene surface,
    \begin{equation}
        \theta(t)=\frac{c_s(t)}{K_d + c_s(t)},
    \end{equation}
    where $\theta(t)$ is fractional site occupancy.
    \item \textbf{Surface charge modulates GFET current:} bound antigen produces an effective electrostatic perturbation that changes the carrier density and results in a measurable drain-current modulation.
    \begin{equation}
        \Delta I_{DS}(t) \propto N(t) \propto \theta(t).
    \end{equation}
\end{enumerate}

Therefore, even before explicitly coupling semiconductor physics, the present study provides the most critical timing insight for a practical sensor: \textbf{heterogeneous diffusion in lymph-node-like tissue creates a measurable exposure delay between compartments.} For an in-body GFET intended to detect HER2, this means sensor placement and transport pathways (e.g., proximity to lymphatic inflow regions) are likely to dominate the early response time and determine how quickly a detectable current shift can be achieved.

\subsection{Full Boundary Exposure}

To investigate how sensor geometry influences detectable signal formation, two sensing configurations were considered. In the first configuration, the entire inlet and outlet boundaries were treated as sensing surfaces, representing an idealized upper-bound scenario where antigen exposure is maximized. 

In the second configuration, only a localized portion of the inlet channel was designated as the sensing region. This setup more closely reflects practical GFET implementations, where the active sensing area is spatially limited.

By comparing these two configurations, the role of transport-driven spatial heterogeneity in determining sensor response can be explicitly evaluated.

% Surface concentration (Gamma) at sensing boundaries
\begin{figure}[H]
    \centering
    \includegraphics[width=0.85\linewidth]{1st surface.png}
    \caption{Surface-bound antigen density ($\Gamma$) distribution at $t = 1000~\text{s}$. The color map represents the local surface concentration of HER2 molecules on selected sensing boundaries. Regions closer to the inlet exhibit higher accumulation, while distal surfaces show significantly lower binding levels, indicating transport-limited exposure.}
    \label{fig:surface_gamma}
\end{figure}

Figure~\ref{fig:surface_gamma} presents the surface-bound antigen density ($\Gamma$) distribution when all inlet and outlet boundaries are treated as sensing regions. This configuration represents an idealized scenario in which the sensor is maximally exposed to incoming analyte.

The results reveal strong spatial heterogeneity in surface accumulation. Surfaces located directly along the primary diffusion pathway exhibit significantly higher antigen binding, whereas regions that are geometrically distant or shielded by low-diffusivity domains show markedly reduced $\Gamma$ values.

This non-uniformity highlights a key limitation: even under conditions of maximal sensing area, antigen exposure remains transport-limited. The availability of analyte at the sensing interface is governed not only by concentration but also by geometric accessibility and diffusion pathways within the domain.

These findings establish an upper bound on achievable sensor response and motivate the investigation of more realistic configurations with spatially localized sensing regions.


\begin{figure}[H]
    \centering
    \includegraphics[width=0.8\linewidth]{plot_c.png}
    \caption{Time evolution of the volume-averaged antigen concentration $c(t)$ under full boundary exposure conditions. The concentration increases monotonically as diffusion distributes the analyte throughout the domain.}
    \label{fig:c_full_boundary}
\end{figure}

Figure~\ref{fig:c_full_boundary} shows the temporal evolution of the volume-averaged antigen concentration within the domain. The concentration increases monotonically as HER2 molecules diffuse from the inlet boundary into the surrounding medium.

The curve exhibits a rapid initial rise followed by a gradual saturation trend, indicating that early-time dynamics are dominated by steep concentration gradients near the inlet. As diffusion progresses, these gradients diminish, leading to a slower rate of concentration increase.

This behavior reflects the transition from a transport-driven regime to a quasi-equilibrium state, where concentration gradients become less pronounced. However, the presence of a continuous increase suggests that complete spatial homogenization is not achieved within the simulated time window.


\begin{figure}[H]
    \centering
    \includegraphics[width=0.8\linewidth]{plot_gamma.png}
    \caption{Time evolution of the average surface-bound antigen density $\Gamma(t)$ for full boundary exposure. The rapid initial increase is followed by a slower accumulation phase, indicating transport-limited binding dynamics.}
    \label{fig:gamma_full_boundary}
\end{figure}

Figure~\ref{fig:gamma_full_boundary} presents the time evolution of the average surface-bound antigen density $\Gamma(t)$ across all sensing boundaries. In contrast to the gradual increase observed in the bulk concentration, the surface response exhibits a sharp initial rise followed by a significantly slower accumulation phase.

This rapid early increase is driven by the immediate availability of antigen molecules near the inlet, allowing fast binding at accessible surfaces. However, as time progresses, the rate of increase in $\Gamma$ becomes limited by the diffusive transport of molecules from deeper regions of the domain.

The logarithmic scale highlights that several orders of magnitude increase occur within a short time window, after which the system enters a transport-limited regime. This behavior indicates that even when the sensing area is maximized, the accumulation of detectable signal is constrained by diffusion rather than surface reaction kinetics.

These results demonstrate that surface binding dynamics do not directly follow bulk concentration trends, and that sensor response is governed by the rate at which analyte molecules reach the sensing interface.

\subsection{Localized Sensing Region}

\begin{figure}[H]
    \centering
    \includegraphics[width=0.75\linewidth]{2nd Surface.png}
    \caption{Localized sensing configuration where only a small region of the inlet channel is designated as the active sensing surface. The color map represents the surface-bound antigen density ($\Gamma$) at $t = 0~\text{s}$, illustrating the spatial restriction of the sensing area.}
    \label{fig:localized_surface_geometry}
\end{figure}

Figure~\ref{fig:localized_surface_geometry} illustrates the localized sensing configuration, in which only a small portion of the inlet channel is treated as the active sensing region. Unlike the full boundary exposure case, this setup restricts antigen detection to a spatially confined surface, better representing realistic GFET implementations where the sensing area is limited by device geometry.

At $t = 0~\text{s}$, the surface-bound antigen density is negligible across the sensing region, as no diffusion-driven transport has yet occurred. This highlights that, in contrast to idealized configurations, signal generation in practical devices depends entirely on the arrival of analyte molecules through transport mechanisms.

The strong spatial localization of the sensing region introduces an additional constraint: only molecules that reach this specific area contribute to the measurable signal. As a result, the effective sensor response is expected to be significantly reduced and delayed compared to the full boundary exposure scenario.



\begin{figure}[H]
    \centering
    \includegraphics[width=0.8\linewidth]{gamma_plot.png}
    \caption{Time evolution of the surface-bound antigen density $\Gamma(t)$ for the localized sensing region. Compared to the full boundary case, significantly lower accumulation is observed due to restricted analyte access.}
    \label{fig:gamma_localized}
\end{figure}

Figure~\ref{fig:gamma_localized} shows the evolution of surface-bound antigen density for the localized sensing configuration. The magnitude of $\Gamma$ is significantly reduced compared to the full boundary exposure case, reflecting the limited accessibility of analyte molecules to the confined sensing region.

Although the temporal trend follows a similar pattern—characterized by a rapid initial increase followed by a slower accumulation phase—the overall signal remains orders of magnitude lower. This confirms that restricting the sensing area introduces a strong transport bottleneck, as only a fraction of diffusing molecules contribute to surface binding.



\begin{figure}[H]
    \centering
    \includegraphics[width=0.8\linewidth]{sigma_plot.png}
    \caption{Effective surface charge density $\sigma_{\text{eff}}(t)$ induced by antigen binding in the localized sensing region. The charge follows the same temporal trend as $\Gamma(t)$, reflecting direct coupling between molecular binding and electrostatic response.}
    \label{fig:sigma_localized}
\end{figure}

Figure~\ref{fig:sigma_localized} presents the effective surface charge density resulting from antigen binding. Since each binding event contributes to local charge accumulation, the temporal behavior of $\sigma_{\text{eff}}(t)$ closely follows that of $\Gamma(t)$.

The reduced magnitude $\sigma_{\text{eff}}$ compared to the full boundary case highlights the direct impact of limited binding events on the electrostatic response. This demonstrates that any reduction in antigen availability at the sensing surface propagates linearly into the electrical domain.

\begin{figure}[H]
    \centering
    \includegraphics[width=0.8\linewidth]{dVg_plot.png}
    \caption{Effective gate voltage shift $\Delta V_g(t)$ induced by surface charge modulation. The response reflects the capacitive coupling between surface charge and the graphene channel.}
    \label{fig:dVg_localized}
\end{figure}

Figure~\ref{fig:dVg_localized} shows the effective gate voltage shift induced by the accumulated surface charge. The relationship between $\sigma_{\text{eff}}$ and $\Delta V_g$ is governed by the electrostatic coupling through the device capacitance.

As expected, the reduced surface charge results in a correspondingly smaller voltage shift. This demonstrates that transport limitations at the molecular level directly translate into reduced modulation of the transistor operating point.


\begin{figure}[H]
    \centering
    \includegraphics[width=0.8\linewidth]{ids_plot.png}
    \caption{Drain-source current shift $\Delta I_{ds}(t)$ resulting from antigen-induced surface charge modulation. This represents the final measurable electrical signal of the GFET biosensor.}
    \label{fig:ids_localized}
\end{figure}

Figure~\ref{fig:ids_localized} illustrates the resulting drain-source current shift, which serves as the measurable output of the GFET sensor. The current response inherits the transport-limited dynamics observed in earlier stages, exhibiting a rapid initial rise followed by a slow approach to saturation.

Importantly, the magnitude $\Delta I_{ds}$ is significantly lower than in the full boundary exposure scenario. This confirms that sensor performance is fundamentally constrained by analyte transport rather than intrinsic device sensitivity.

The results across Figures~\ref{fig:gamma_localized}–\ref{fig:ids_localized} demonstrate a clear propagation of transport limitations through the entire sensing chain. Restricted analyte access reduces surface binding, which in turn limits surface charge accumulation, gate voltage modulation, and ultimately the measurable current response.

This cascading attenuation highlights that the overall performance of GFET-based biosensors is not solely determined by electronic sensitivity but is fundamentally governed by the efficiency of analyte transport to the sensing interface.

\subsection{Comparison and Engineering Implications}

The results obtained from the full boundary exposure and localized sensing configurations reveal a fundamental difference in sensor behavior arising from geometric constraints.

In the full boundary case, the sensing area spans all inlet and outlet surfaces, providing maximal exposure to the incoming analyte. This configuration yields relatively high surface binding density and correspondingly strong electrical signals. However, such a setup represents an idealized upper bound and does not reflect realistic device geometries.

In contrast, the localized sensing configuration restricts detection to a small region of the inlet channel, significantly limiting analyte accessibility. As a result, the surface binding density is reduced by several orders of magnitude, and this reduction propagates through the entire signal transduction chain. The effective surface charge, gate voltage shift, and drain current response all exhibit similarly attenuated behavior.

A key observation is that the temporal dynamics of both configurations follow a similar pattern: a rapid initial increase followed by a slower accumulation phase. However, in the localized case, the transition to the transport-limited regime occurs earlier and persists for a longer duration. This indicates that diffusion constraints dominate the sensing process much more strongly when the active sensing area is spatially confined.

These findings highlight a critical design principle for GFET-based biosensors: sensor performance is fundamentally governed by analyte transport rather than intrinsic electronic sensitivity. Even with highly responsive transducer elements, the measurable signal remains limited if analyte molecules cannot efficiently reach the sensing interface.

From an engineering perspective, this implies that improving device performance requires not only optimizing transistor characteristics but also carefully designing the spatial configuration of the sensing region. Strategies such as increasing effective sensing area, positioning sensors along dominant transport pathways, or enhancing local transport conditions may be necessary to achieve meaningful signal levels.

Overall, the results demonstrate that realistic biosensor operation must be analyzed as a coupled transport - reaction - electronic system, where limitations at the transport level directly constrain the final electrical output.



\appendix
\chapter{Appendix}
\section{Key Equations}
\begin{align}
  \Delta I_{ds} &= \frac{w}{l} e \mu V_{ds} \Delta n \\
  \theta &= \frac{C}{C + K_d} \\
  LOD &= \frac{3\sigma}{S}
\end{align}

\section{Simulation Parameters}
\begin{center}
\begin{tabular}{|l|c|c|}
\hline
\textbf{Parameter} & \textbf{Value} & \textbf{Unit} \\
\hline
Channel Length (L) & 5 & $\mu$m \\
Channel Width (W) & 20 & $\mu$m \\
Carrier Mobility ($\mu$) & 2000 & cm$^2$/V$\cdot$s \\
Source-Drain Voltage ($V_{ds}$) & 0.1 & V \\
Gate Bias ($V_g^*$) & 0.15 & V \\
Debye Length & 1-4 & nm \\
HER2 Concentration Range & 0.5-1000 & pM \\
\hline
\end{tabular}
\end{center}

\section{Expected Outcomes}
\begin{itemize}
  \item Approximate linear correlation between bound HER2 molecules and $\Delta I_{ds}$ under low-coverage conditions; deviations may occur due to transport limitations.
  \item Optimum gate bias minimizing noise and maximizing sensitivity.
  \item Predicted limit of detection around 0.5-1.0 pM under low-salt conditions.
\end{itemize}

\section{Node-Module Parameter Summary}

\begin{table}[H]
\centering
\begin{tabular}{lcc}
\toprule
Parameter & Symbol / Value & Unit \\ \midrule
Channel length & $l = 2~\mu$m & -- \\
Channel width & $w = 20~\mu$m & -- \\
Carrier mobility & $\mu = 0.1$ & m$^2$/V·s \\
Drain bias & $V_{ds} = 50$ & mV \\
Effective area & $A_{\text{eff}} = 4\times10^{-11}$ & m$^2$ \\
Coupling efficiency & $\alpha = 0.01$--$0.03$ & -- \\
Current noise floor & $\Delta I_{ds,\min} = 10$--$50$ & pA \\
Minimum bound number & $N_{\min} \approx 5$--$25$ & molecules \\ \bottomrule
\end{tabular}
\caption{Realistic node-module GFET parameters and derived sensitivity range.}
\end{table}

\noindent
These parameters define a practical lower bound for in-node HER2 detection:
tens of bound molecules yield a resolvable current shift, consistent with the
picomolar LOD obtained in simulation.

\section{Transport Model Definition}

Antigen transport was modeled using the \textit{Transport of Diluted Species} interface in COMSOL Multiphysics under diffusion-only assumptions. Convection and electromigration terms were explicitly disabled to isolate the effect of heterogeneous diffusion.

The governing equation solved over the three-dimensional domain was:
\begin{equation}
\frac{\partial c}{\partial t}
= \nabla \cdot \left( D(\mathbf{x}) \nabla c \right),
\end{equation}
where $c(\mathbf{x},t)$ denotes the HER2 concentration and $D(\mathbf{x})$ is a piecewise-constant diffusion coefficient assigned per subdomain.

This formulation isolates transport limitations, enabling direct assessment of how analyte delivery constrains downstream signal generation in the GFET sensing chain.

\section{Signal Transduction Model}

To connect transport dynamics with measurable electrical response, a multi-stage signal transduction model was employed. The process links molecular binding events to the final drain-source current shift of the GFET device.

First, antigen transport determines the local concentration at the sensing surface, which governs the surface binding density $\Gamma(t)$. The number of bound molecules is then related to surface charge density through:
\begin{equation}
\sigma_{\text{eff}}(t) = q \cdot \Gamma(t),
\end{equation}
where $q$ represents the effective charge contribution per bound molecule.

The accumulated surface charge induces a shift in the effective gate voltage through capacitive coupling:
\begin{equation}
\Delta V_g(t) = \frac{\sigma_{\text{eff}}(t)}{C_{\text{eq}}},
\end{equation}
where $C_{\text{eq}}$ is the equivalent capacitance of the sensing interface, including contributions from the electrical double layer and graphene channel.

Finally, the gate voltage modulation alters the carrier density in the graphene channel, producing a measurable drain-source current shift:
\begin{equation}
\Delta I_{ds}(t) = \frac{w}{l} e \mu V_{ds} \Delta n(t),
\end{equation}
where $\Delta n(t)$ is directly proportional to $\Delta V_g(t)$.

This formulation establishes a direct causal chain:
\[
\text{Transport} \rightarrow \Gamma(t) \rightarrow \sigma(t) \rightarrow \Delta V_g(t) \rightarrow \Delta I_{ds}(t),
\]
allowing the impact of transport limitations to be traced through all stages of signal formation.

\section{Domain Partitioning and Volume Selections}

The computational domain was partitioned into two biologically motivated regions:
\begin{itemize}
    \item $\Omega_{\text{cortex}}$: outer shell representing the lymph node cortex and capsule,
    \item $\Omega_{\text{medulla}}$: inner volume representing the medulla.
\end{itemize}

Explicit domain selections were defined in COMSOL and used consistently for:
\begin{itemize}
    \item Assigning heterogeneous diffusion coefficients,
    \item Computing volume-averaged concentrations,
    \item Extracting compartment-specific uptake curves.
\end{itemize}

Volume averages were computed using:
\begin{equation}
\langle c(t) \rangle_{\Omega_i}
=
\frac{1}{|\Omega_i|}
\int_{\Omega_i} c(\mathbf{x},t)\, dV,
\qquad
\Omega_i \in \{\Omega_{\text{cortex}}, \Omega_{\text{medulla}}\}.
\end{equation}

\section{Diffusion Parameters and Parametric Sweep}

Diffusion coefficients were assigned independently to cortex and medulla subdomains to represent transport heterogeneity:
\begin{align}
D_{\text{cortex}} &= D_0, \\
D_{\text{medulla}} &= r \cdot D_0,
\end{align}
where $r$ is the diffusivity ratio.

A parametric sweep was performed over:
\[
r \in \{1.0, 0.7, 0.5, 0.3, 0.2, 0.1\}.
\]

This sweep quantifies how increasingly restrictive internal transport delays antigen exposure in the medulla relative to the cortex.


\section{Simulation Time Horizon and Solver Settings}

Transient simulations were conducted over two characteristic time windows:
\begin{itemize}
    \item Short-term exposure: $t \leq 200~\text{s}$, used to evaluate early transport and inlet-dominated gradients.
    \item Long-term exposure: $t \leq 6000~\text{s}$, used to assess quasi-steady diffusion and deep tissue penetration.
\end{itemize}

The extended time horizon ensures that observed cortex-medulla delays are not artifacts of insufficient simulation duration. All simulations used adaptive time stepping with solver tolerances selected to ensure numerical stability and mass conservation.

\section{Assumptions and Limitations}

\begin{itemize}
    \item No convection or interstitial flow was considered; transport is purely diffusive.
    \item Binding kinetics were assumed to be transport-limited, neglecting reaction-rate limitations.
    \item Electrostatic coupling was modeled using a linear capacitance approximation.
    \item Debye screening effects were incorporated implicitly through effective parameters.
\end{itemize}

 

\begin{thebibliography}{9}
\addcontentsline{toc}{chapter}{Bibliography}
\bibitem{janeway2001} C.~A.~Janeway \textit{et al.}, \textit{Immunobiology: The Immune System in Health and Disease}, 5th~ed. New York, NY: Garland Science, 2001.
\bibitem{weissleder2004} M.~R.~Weissleder and U.~Pittet, ``Imaging of lymphatic system and metastasis,'' \textit{N. Engl. J. Med.}, vol.~351, no.~18, pp.~1800-1802, 2004.
\bibitem{yarden2001} S.~Yarden and M.~X.~Sliwkowski, ``Untangling the ErbB signalling network,'' \textit{Nat. Rev. Mol. Cell Biol.}, vol.~2, pp.~127-137, 2001.
\bibitem{kabel2017} A.~M.~Kabel, ``Tumor markers of breast cancer: New prospectives,'' \textit{J. Oncol. Sci.}, vol.~3, no.~1, pp.~5-11, Apr.~2017.
\bibitem{harpaz2017} D.~Harpaz \textit{et al.}, ``Point-of-care-testing in acute stroke management: An unmet need ripe for technological harvest,'' \textit{Biosensors (Basel)}, vol.~7, no.~3, p.~30, 2017.
\bibitem{janeway2001ch9} C.~A.~Janeway \textit{et al.}, \textit{Immunobiology}, ch.~9, pp.~280-285, 2001.
\bibitem{bergveld1972} P.~Bergveld, ``Development, operation, and application of the ion-sensitive field-effect transistor as a tool for electrophysiology,'' \textit{IEEE Trans. Biomed. Eng.}, vol.~BME-19, no.~5, pp.~342-351, Sep.~1972.
\bibitem{pena2018} J.~Peña-Bahamonde \textit{et al.}, ``Recent advances in graphene-based biosensor technology with applications in life sciences,'' \textit{J. Nanobiotechnol.}, vol.~16, no.~1, p.~75, 2018.
\bibitem{biodigital}BioDigital, Inc., ``BioDigital Human: 3D Human Visualization Platform,'' 2025. [Online]. Available: \url{https://human.biodigital.com/}
\bibitem{demoraes2016} A.~C.~De Moraes and L.~T.~Kubota, ``Recent trends in field-effect transistors-based immunosensors,'' \textit{Chemosensors}, vol.~4, no.~4, 2016.
\bibitem{seo2020} G.~Seo \textit{et al.}, ``Rapid detection of COVID-19 causative virus (SARS-CoV-2) in human nasopharyngeal swab specimens using field-effect transistor-based biosensor,'' \textit{ACS Nano}, vol.~14, no.~4, pp.~5135-5142, Apr.~2020.
\bibitem{kim2018} T.~Kim, M.~Cho, and K.~J.~Yu, ``Flexible and stretchable bio-integrated electronics based on carbon nanotube and graphene,'' \textit{Materials}, vol.~11, no.~7, 2018.
\bibitem{xu2022} S.~Xu \textit{et al.}, ``Highly sensitive HER2 biosensor based on graphene oxide field-effect transistor for breast cancer diagnosis,'' \textit{Biosens. Bioelectron.}, vol.~195, p.~113691, 2022.
\bibitem{zhang2021} Z.~Zhang \textit{et al.}, ``Multiplexed detection of cancer biomarkers using graphene field-effect transistor biosensors,'' \textit{ACS Appl. Mater. Interfaces}, vol.~13, no.~35, pp.~41546-41556, 2021.
\bibitem{iso10993}
International Organization for Standardization, ``ISO 10993-1: Biological evaluation of medical devices – Part 1: Evaluation and testing within a risk management process,'' ISO, Geneva, Switzerland, 2018.

\end{thebibliography}

\chapter*{Acknowledgments}
\addcontentsline{toc}{chapter}{Acknowledgments}


Sincere thanks are extended to \textbf{M. Görkem Durmaz}, whose initial BAP project on graphene-based biosensor modeling provided the conceptual and scientific groundwork for this investigation. Her thorough preparation and documentation gave this report crucial theoretical understanding and methodological guidance.  


\end{document}


